% This is samplepaper.tex, a sample chapter demonstrating the
% LLNCS macro package for Springer Computer Science proceedings;
% Version 2.21 of 2022/01/12
%
\documentclass[runningheads]{llncs}
%
\usepackage[T1]{fontenc}
% T1 fonts will be used to generate the final print and online PDFs,
% so please use T1 fonts in your manuscript whenever possible.
% Other font encondings may result in incorrect characters.
%
\usepackage{graphicx}
\usepackage{amsmath,amssymb}
\usepackage{hyperref}
% Used for displaying a sample figure. If possible, figure files should
% be included in EPS format.
%
% If you use the hyperref package, please uncomment the following two lines
% to display URLs in blue roman font according to Springer's eBook style:
%\usepackage{color}
%\renewcommand\UrlFont{\color{blue}\rmfamily}
%\urlstyle{rm}
%
\begin{document}
%
\title{Adaptive Risk-aware Implicit Quantile Network: Towards Safe and Energy-Efficient USV Navigation in Dynamic Environments}
%
%\titlerunning{Abbreviated paper title}
% If the paper title is too long for the running head, you can set
% an abbreviated paper title here
%
% \author{First Author\inst{1}\orcidID{0000-1111-2222-3333} \and
% Second Author\inst{2,3}\orcidID{1111-2222-3333-4444} \and
% Third Author\inst{3}\orcidID{2222--3333-4444-5555}}
% %
% \authorrunning{F. Author et al.}
% % First names are abbreviated in the running head.
% % If there are more than two authors, 'et al.' is used.
% %
% \institute{Princeton University, Princeton NJ 08544, USA \and
% Springer Heidelberg, Tiergartenstr. 17, 69121 Heidelberg, Germany
% \email{lncs@springer.com}\\
% \url{http://www.springer.com/gp/computer-science/lncs} \and
% ABC Institute, Rupert-Karls-University Heidelberg, Heidelberg, Germany\\
% \email{\{abc,lncs\}@uni-heidelberg.de}}
% %
\maketitle              % typeset the header of the contribution
\begin{abstract}
Safe and efficient local navigation for unmanned surface vehicles (USVs) remains challenging in unknown, time-varying current fields with dynamic obstacles. Conventional deep reinforcement learning methods typically optimize expected return and may therefore underestimate rare but catastrophic failures caused by current--obstacle interactions. To address this issue, we propose \textbf{Adaptive Risk-aware Implicit Quantile Network (Adaptive Risk-aware IQN)}, a distributional reinforcement learning framework for sensor-driven USV navigation. Our method models the full return distribution with IQN and applies a CVaR-style risk distortion to emphasize adverse tail outcomes during action selection. To avoid excessive conservatism, we further design a distance-adaptive risk mechanism that increases risk aversion near obstacles while remaining nearly risk-neutral in open water. Experiments in dynamic marine environments show that our method achieves a superior balance between safety and efficiency, reaching a \textbf{97\% success rate} in the hardest setting while maintaining strong time and energy performance.
\keywords{Unmanned surface vehicle, distributional reinforcement learning, implicit quantile network, conditional value-at-risk, dynamic environment}
\end{abstract}
%
%

\section{Introduction}

Autonomous surface vehicles (USVs) operating in ocean and riverine environments must cope with complex flow disturbances, where achieving both safety and efficiency remains a fundamental challenge~\cite{ref1}. 
With the increasing deployment of USVs in real-world scenarios such as environmental monitoring, maritime operations, and field competitions, reliable navigation under uncertain dynamics has become a critical requirement. 
In particular, local navigation in unknown and time-varying current fields poses significant challenges, as the vehicle must make decisions based solely on onboard perception without prior knowledge of the environment~\cite{ref5}. 
This difficulty is further exacerbated by the presence of dynamic obstacles, which introduces highly stochastic interactions and increases the likelihood of rare but catastrophic failures.

Reinforcement learning (RL) provides a promising paradigm for learning navigation policies from interaction data, enabling generalization to previously unseen environments without requiring explicit modeling of complex dynamics~\cite{ref6}. 
In particular, deep reinforcement learning (DRL) has demonstrated strong performance in high-dimensional decision-making tasks, including robot navigation. 
However, most existing DRL approaches optimize the \emph{expected return}, implicitly assuming that average performance is sufficient. 
In safety-critical navigation scenarios, this assumption can lead to unreliable behavior, as rare but high-impact failures—such as collisions caused by flow--obstacle interactions—are largely ignored under expectation-based optimization.

To better capture such uncertainty, distributional reinforcement learning has been proposed to model the \emph{full distribution} of returns rather than only their expectation~\cite{ref7}. 
By representing the variability of outcomes, distributional methods enable risk-sensitive decision-making through techniques such as risk distortion~\cite{ref8}. 
Nevertheless, existing approaches typically rely on fixed risk preferences, which fail to adapt to rapidly changing environments. 
As a result, the learned policies may become overly conservative in safe regions or insufficiently cautious near hazards, leading to suboptimal trade-offs between safety and efficiency.
\begin{figure}[htbp]
  \centering
  \includegraphics[width=0.85\textwidth]{figures/intro.png}
  \caption{Navigation in uncertain flow fields with dynamic obstacles. Other methods either fail due to risk underestimation or are inefficient. Our Adaptive IQN proactively avoids hazards and reliably reaches the target.}
  \label{fig:intro_overview}
\end{figure}
To address these challenges, we propose \textbf{Adaptive Risk-Aware IQN}, a distributional reinforcement learning framework for sensor-driven USV navigation under unknown current fields with dynamic obstacles. 
The proposed framework models return distributions using Implicit Quantile Networks (IQN) and performs risk-aware action selection via a distortion mechanism that emphasizes adverse outcomes. 
Furthermore, we introduce an adaptive risk modulation strategy that dynamically adjusts risk sensitivity according to environmental conditions, encouraging conservative behavior near hazards while maintaining efficient navigation in safer regions. 
In addition, we incorporate an energy-aware objective into the learning process, enabling the agent to achieve safer navigation with reduced control effort. 
Fig.~\ref{fig:intro_overview} illustrates a representative scenario, where conventional methods either fail due to risk underestimation or suffer from inefficiency, while our approach successfully achieves safe and efficient navigation.

\smallskip
\noindent The main contributions of this work are as follows:

\begin{enumerate}

\item \textbf{A risk-aware distributional reinforcement learning framework for USV navigation.}
We propose an Adaptive Risk-Aware IQN framework that explicitly models return distributions and enables risk-sensitive decision-making, significantly improving robustness and navigation success rate in dynamic environments with unknown currents and moving obstacles.

\item \textbf{An adaptive risk modulation mechanism.}
We design a dynamic strategy that adjusts risk sensitivity online based on environmental context, allowing the agent to behave conservatively near hazards while maintaining efficient navigation in open regions.

\item \textbf{An energy-aware navigation objective.}
We introduce an energy optimization term into the learning objective, enabling the agent to reduce control effort while preserving strong safety performance, which is critical for long-duration deployment.

\item \textbf{Extensive evaluation in dynamic flow environments.}
We conduct comprehensive experiments and compare against classical planners and strong DRL baselines. Results demonstrate that our method consistently achieves superior performance in terms of safety, efficiency, and energy consumption.

\end{enumerate}
\section{Related Work}
Classical reactive planners, learning-based navigation methods, and risk-sensitive Distributional RL have been widely studied for mobile robot and marine vehicle navigation. In our setting, the challenge is further amplified by the joint effect of dynamic obstacles and unknown, spatially varying current fields, which creates highly nonstationary and safety-critical local interactions.

Artificial Potential Field (APF) methods are classical reactive planners known for their simplicity and real-time efficiency~\cite{ref13}. A number of later studies have improved APF to alleviate local minima, deadlock, or poor behavior in dynamic environments~\cite{ref16,ref17}. Bug-family algorithms provide another lightweight navigation paradigm based on local sensing and obstacle-boundary following~\cite{ref22,ref23}. Subsequent variants further improve path efficiency and practical performance in robotic navigation~\cite{ref28,ref30}. Despite their simplicity and online responsiveness, these reactive planners do not explicitly model long-horizon uncertainty or safety-critical tail events. Their performance can therefore degrade when moving obstacles and external current disturbances jointly create rapidly changing hazard regions.

Learning-based navigation methods, especially Deep Reinforcement Learning (DRL), improve adaptability by learning policies directly from interaction data. Among value-based methods, DQN and its extensions, including DDQN, dueling networks, and prioritized experience replay, improve value estimation stability and sample efficiency~\cite{ref39,ref40,ref41,ref42}. These methods have also been applied to navigation and control tasks in challenging environments~\cite{ref31,ref32,ref34,ref35}. Their combination, often referred to in practice as D3QN, has also been used in dynamic path planning~\cite{ref43}. In addition, policy-gradient and actor--critic methods such as PPO have been widely adopted for navigation and continuous-control tasks because of their stable on-policy updates and practical robustness~\cite{ref44,ref45}. Compared with hand-crafted reactive planners, these methods better capture sequential decision structure and adapt behavior from experience. Nevertheless, most of them remain \emph{risk-neutral} in the sense that they optimize the expected return. In safety-critical navigation, this can be problematic because rare but catastrophic collision events may contribute little to the expectation while still dominating practical risk.

This limitation has motivated increasing interest in Distributional RL, which models the full return distribution rather than only its expectation~\cite{ref7,ref8}. In particular, IQN represent the value distribution through sampled quantile levels, providing a natural mechanism for risk-sensitive decision making by emphasizing different parts of the return distribution during action selection~\cite{ref8}. Recent studies suggest that distributional methods can improve safety and robustness in navigation tasks~\cite{ref36,ref37}, and some work has further explored adaptive risk modulation according to environmental difficulty~\cite{ref38}. However, prior work has only partially addressed the setting considered here: the coupled effect of unknown current fields and dynamic obstacles remains underexplored, and existing approaches often lack online adaptation of risk sensitivity to local hazard conditions. Our work builds on this line by combining IQN-style distributional value learning with adaptive risk modulation for USV navigation in unknown current-disturbed environments with dynamic obstacles.


\section{Problem Formulation}
We model the USV navigation task as a Markov Decision Process (MDP) defined by the tuple $(\mathcal{S}, \mathcal{A}, P, R, \gamma)$, where $\mathcal{S}$ and $\mathcal{A}$ denote the state space and action space, respectively. At each discrete time step $t$, the USV observes the current state $s_t \in \mathcal{S}$ and chooses an action $a_t \in \mathcal{A}$. The system then evolves to the next state according to the transition dynamics $s_{t+1} \sim P(\cdot \mid s_t, a_t)$, after which the agent receives a scalar reward $r_{t+1} = R(s_t, a_t)$ and proceeds with the subsequent decision.

A standard objective in reinforcement learning is to learn a policy $\pi$ that maximizes the expected discounted cumulative return. This objective is captured by the action-value function (or $Q$-function) under policy $\pi$, defined as
\begin{equation}
Q^{\pi}(s,a) = \mathbb{E}_{\pi}\!\left[\sum_{k=0}^{\infty}\gamma^{k}\, r_{t+k+1}\,\Big|\, s_t=s,\, a_t=a \right],
\label{eq:q_def}
\end{equation}
where $\gamma \in [0,1)$ is the discount factor controlling the relative importance of future rewards. Starting from a given state--action pair $(s_t,a_t)$, the expectation is taken over future transitions and actions when following $\pi$ thereafter.

The $Q$-function satisfies the Bellman expectation equation
\begin{equation}
Q^{\pi}(s,a) = \mathbb{E}\!\left[ R(s,a) \right] + \gamma\, \mathbb{E}\!\left[ Q^{\pi}(s',a') \right],
\label{eq:bellman_exp}
\end{equation}
where $s' \sim P(\cdot \mid s,a)$ and $a' \sim \pi(\cdot \mid s')$. In particular, an optimal policy can be obtained by applying the Bellman optimality operator, which is commonly written as
\begin{equation}
\mathcal{T}Q(s,a) := \mathbb{E}\!\left[ R(s,a) \right] + \gamma\, \mathbb{E}\!\left[\max_{a'} Q(s',a') \right],
\label{eq:bellman_opt}
\end{equation}
with $s' \sim P(\cdot \mid s,a)$. The fixed point of this operator corresponds to the optimal action-value function, from which a greedy policy can be derived.


\section{Methods}

\label{sec:method}

\subsection{Traditional DRL Baseline}
\label{sec:method_dqn}
We adopt a value-based deep reinforcement learning baseline to represent conventional scalar-value training. Specifically, we focus on the foundational Deep Q-Network (DQN) mechanics to learn an action-value function $Q(s,a;\theta)$, which approximates the expected discounted return defined in Eq.~(1). Given transition samples $(s,a,r,s')$ collected from interaction, it minimizes a temporal-difference (TD) regression objective with a target network $\theta^-$:
\begin{equation}
\mathcal{L}_{\mathrm{DQN}}
=
\mathbb{E}\Big[
\big(
r + \gamma \max_{a'} Q(s',a';\theta^-)
- Q(s,a;\theta)
\big)^2
\Big].
\label{eq:dqn_loss}
\end{equation}
In our experiments, we employ D3QN---an advanced combination of DQN, Double DQN, and Dueling architectures---as a representative risk-neutral baseline to provide a stronger, more comprehensive comparison against our method.

\subsection{Distributional RL}
\label{sec:method_distRL}
Unlike DQN, which learns only the expectation of returns, distributional reinforcement learning models the \emph{return distribution} for each state--action pair. Let $Z^{\pi}(s,a)$ denote the random return under policy $\pi$, satisfying $\mathbb{E}[Z^{\pi}(s,a)] = Q^{\pi}(s,a)$. The distributional Bellman relation can be written as
\begin{equation}
Z^{\pi}(s,a) \overset{D}{=} R(s,a) + \gamma Z^{\pi}(s',a'),
\label{eq:dist_bellman}
\end{equation}
where $s' \sim P(\cdot|s,a)$ and $a' \sim \pi(\cdot|s')$, and $\overset{D}{=}$ denotes equality in distribution.
To obtain an optimal distributional value function, we apply the distributional Bellman optimality operator, where the greedy action is determined by the expectation of the return distribution:
\begin{equation}
\mathcal{T}Z(s,a) \overset{D}{=} R(s,a) + \gamma Z\!\Big(s',\,\arg\max_{a'} \mathbb{E}[Z(s',a')]\Big).
\label{eq:dist_opt}
\end{equation}

Our algorithm employs IQN to represent $Z(s,a)$ via its quantile function. As shown in Fig.~\ref{fig:overview_arch}, IQN predicts action-conditioned return \emph{quantiles} by combining a fused state feature with an embedding of sampled quantile fractions.
For $\tau \sim \mathcal{U}([0,1])$, define the $\tau$-quantile as
\begin{equation}
Z_{\tau}(s,a) := F_Z^{-1}(\tau),
\label{eq:iqn_quantile}
\end{equation}
where $F_Z$ is the CDF of the return distribution. Risk sensitivity is introduced through a distortion function $\beta:[0,1]\rightarrow[0,1]$, yielding a distorted expectation
\begin{equation}
Q_{\beta}(s,a)
=
\mathbb{E}_{\tau \sim \mathcal{U}([0,1])}\Big[ Z_{\beta(\tau)}(s,a)\Big].
\label{eq:distorted_q}
\end{equation}

% -----------------------------
% (1) Added: explicit CVaR-style distortion definition
% -----------------------------
In particular, We adopt a CVaR-style distortion parameterized by $\phi\in(0,1]$, namely $\beta_{\phi}(\tau)=\phi\tau$ for $\tau\in[0,1]$,
which corresponds to emphasizing the lower-tail outcomes by averaging the lowest $\phi$ fraction of quantiles (i.e., a CVaR-type criterion). Intuitively, $\beta(\cdot)$ reshapes the contribution of quantiles so that unfavorable outcomes exert greater influence on the decision process. As illustrated in Fig.~\ref{fig:overview_arch}, the distorted quantiles are subsequently aggregated into a risk-aware value used for action selection.

The resulting risk-aware policy is
\begin{equation}
\pi_{\beta}(s)=\arg\max_{a} Q_{\beta}(s,a),
\label{eq:risk_policy}
\end{equation}
which we approximate by Monte Carlo sampling with $K$ quantiles:
\begin{equation}
\hat{\pi}_{\beta}(s)=
\arg\max_{a}\frac{1}{K}\sum_{k=1}^{K} Z_{\beta(\tilde{\tau}_k)}(s,a),
\quad
\tilde{\tau}_k \sim \mathcal{U}([0,1]).
\label{eq:risk_policy_mc}
\end{equation}

IQN is trained using a quantile regression formulation. Following the standard IQN training procedure, given sampled quantiles $\{\tau_i\}_{i=1}^{N}$ and $\{\tau'_j\}_{j=1}^{N'}$, we form the sampled TD error
\begin{equation}
\delta_{\tau_i,\tau'_j}
=
r + \gamma Z_{\tau'_j}\!\big(s', \hat{\pi}_{\beta}(s')\big)
- Z_{\tau_i}(s,a),
\label{eq:iqn_td}
\end{equation}
and minimize the quantile Huber loss:
\begin{align}
\rho^{\kappa}_{\tau}(u)
&=
|\tau-\mathbb{I}\{u<0\}|\cdot \frac{\mathcal{L}_{\kappa}(u)}{\kappa},
\label{eq:qhuber}\\
\mathcal{L}_{\kappa}(u)
&=
\begin{cases}
\frac{1}{2}u^{2}, & |u|\le\kappa,\\
\kappa(|u|-\frac{1}{2}\kappa), & \text{otherwise}.
\end{cases}
\nonumber
\end{align}
The final loss is
\begin{equation}
\mathcal{L}_{\mathrm{IQN}}
=
\frac{1}{N'}\sum_{i=1}^{N}\sum_{j=1}^{N'}
\rho^{\kappa}_{\tau_i}\!\left(\delta_{\tau_i,\tau'_j}\right).
\label{eq:iqn_loss}
\end{equation}

\subsection{States and Observations}
\label{sec:method_obs}
The full environment state contains the USV pose and motion variables, the (unknown) current field, the goal location, and obstacle states (including dynamic obstacles). In practice, the model operates under partial observability and receives onboard measurements composed of three modalities: $o_t=\big(o^{\mathrm{vel}}_t,\; o^{\mathrm{goal}}_t,\; o^{\mathrm{lidar}}_t\big)$.
These modalities are encoded by separate MLPs and fused via concatenation to form the model input feature.
Following the sensor assumptions in the reference setting, $o^{\mathrm{vel}}_t$ represents the seafloor-relative velocity measured by a Doppler Velocity Log (DVL), $o^{\mathrm{goal}}_t$ denotes the goal position expressed in the robot frame using navigation sensors (e.g., IMU/compass/GPS), and $o^{\mathrm{lidar}}_t$ encodes range reflections indicating nearby obstacles. We use a fixed LiDAR sensing range $d_0$ and treat both static and dynamic obstacles uniformly through their instantaneous detections. Dynamic obstacles additionally change position over time, increasing uncertainty and making risk-aware decision making more critical.

\subsection{Navigation Environment}
\label{sec:method_env}
We construct a simulated marine navigation environment in which the USV must reach a goal under flow disturbances while avoiding obstacles. The current field is generated using a superposition of vortices, where each vortex follows a Rankine model with core radius $r_0$ and circulation strength $\Gamma$. In polar coordinates centered at the vortex, the flow satisfies
\begin{equation}
v_r = 0,\quad
v_{\theta}(r)=\frac{\Gamma}{2\pi}
\begin{cases}
\frac{r}{r_0^2}, & r \le r_0,\\
\frac{1}{r}, & r > r_0.
\end{cases}
\label{eq:rankine}
\end{equation}
The net flow velocity at a location is obtained by summing contributions from nearby vortices. The USV is modeled kinematically as a point particle whose total velocity is the sum of the ambient current velocity and the commanded steering velocity:
\begin{equation}
\frac{dX(t)}{dt}=V(t)=V_C(X(t)) + V_S(t),
\label{eq:kinematics}
\end{equation}
where $X(t)$ is the USV position, $V_C(\cdot)$ is the local current velocity, and $V_S(t)$ is the controllable steering velocity (cf. the control process described above).
Episodes terminate upon reaching the goal, colliding with an obstacle, leaving the workspace boundary, or exceeding a maximum horizon.

\subsection{Action Space}
\label{sec:method_action}

We discretize the control command as a change in the steering velocity $V_S$. Concretely, at each time step the agent selects an action $a_t=(a,w)$.
% -----------------------------
% (2) Added: discrete-time update for V_S from (a,w)
% -----------------------------
We implement the kinematics in discrete time with control interval $\Delta t$. Let $v_t=\|V_S(t)\|$ and $\theta_t$ denote the magnitude and direction of the steering velocity, respectively. The action $(a,w)$ updates the steering command as
\begin{align}
v_{t+1} &= \mathrm{clip}\!\left(v_t + a\,\Delta t,\; 0,\; v_{\max}\right), \label{eq:vs_update_v}\\
\theta_{t+1} &= \theta_t + w\,\Delta t, \label{eq:vs_update_theta}\\
V_S(t+1) &= v_{t+1}\big[\cos(\theta_{t+1}),\,\sin(\theta_{t+1})\big]. \label{eq:vs_update_vec}
\end{align}

where $a$ adjusts the magnitude of $V_S$ (interpretable as a linear acceleration command) and $w$ adjusts its direction (an angular velocity command). We use a small discrete set for each component (e.g., $a \in \{-a_0,0,a_0\}$ and $w \in \{-w_0,0,w_0\}$), and clip the resulting forward speed to $[0,v_{\max}]$. This discrete parameterization enables direct comparison between D3QN and IQN under a shared action interface.

\subsection{Reward Function}
\label{sec:method_reward}
The reward is designed to encourage steady progress toward the goal while penalizing collisions and excessive energy usage. The agent receives a scalar reward after executing each action, where the reward incorporates both task progress and an energy-related penalty. We use a shaped step reward based on goal-distance reduction and terminal bonuses/penalties:
\begin{equation}
r_t=
\begin{cases}
r_{\mathrm{base},t} + r_{\mathrm{col}}, & \text{if collision at } t,\\
r_{\mathrm{base},t} + r_{\mathrm{goal}}, & \text{if goal reached at } t,\\
r_{\mathrm{base},t}, & \text{otherwise},
\end{cases}
\label{eq:reward_piecewise}
\end{equation}
with
\begin{equation}
r_{\mathrm{base},t}
=
r_{\mathrm{step}} + \alpha\big(d_{t-1}-d_t\big) - \lambda_E\,\mathcal{E}(a_t).
\label{eq:reward_base}
\end{equation}
Here $d_t$ is the Euclidean distance from the USV to the goal at time $t$, and $r_{\mathrm{step}}<0$ encourages efficiency. The term $\mathcal{E}(a_t)$ is an energy-related penalty (e.g., proportional to control effort or commanded speed change), and $\lambda_E$ controls the energy--performance trade-off. This explicit energy term supports our energy consumption analysis and encourages trajectories that are not only safe but also energy efficient.

To explicitly quantify the \emph{propulsion work} contributed by the USV itself, we define $\mathcal{E}(a_t)$ using the commanded steering velocity $V_S(t)$ in Eq.~\eqref{eq:kinematics}. Specifically, let $\Delta t$ be the control interval and define the self-propulsion displacement as $\Delta X_{S,t}=V_S(t)\Delta t$. We then set
\begin{equation}
\mathcal{E}(a_t) = \|\Delta X_{S,t}\| = \|V_S(t)\|\Delta t,
\label{eq:energy_def}
\end{equation}
which serves as a simple energy proxy proportional to the magnitude of the control-induced motion. With this definition, $\mathcal{E}(a_t)=0$ when the vehicle drifts with the current (i.e., $V_S(t)=0$), aligning with our goal of encouraging energy-efficient navigation that leverages favorable currents whenever possible.

\begin{figure*}[t]
  \centering
  \includegraphics[width=\textwidth]{figures/f1.pdf}
  \caption{The Overall Framework of Adaptive Risk-aware IQN.}
  \label{fig:overview_arch}
\end{figure*}

\subsection{Algorithm Architecture}
\label{sec:method_arch}
As shown in Fig.~\ref{fig:overview_arch}, the proposed model first encodes velocity, goal-relative measurements, and LiDAR observations with separate multilayer perceptrons, and concatenates them into a fused state representation. For comparison, the standard scalar-value baseline uses the same observation encoder but replaces the IQN distribution head with standard scalar action-value prediction.

For IQN, in addition to the fused state feature, we sample quantile fractions $\tau$ and embed them using a cosine feature map. The quantile embedding is combined with the state feature via element-wise multiplication, followed by fully connected layers to output a set of action-conditioned quantile values $\{Z_{\tau}(s,a)\}$. During decision making, the adaptive risk module maps the obstacle distance $d_{\min}$ to a risk sensitivity parameter $\phi$, applies the distortion $\beta_{\phi}(\tau)=\phi\tau$, and aggregates the distorted quantiles into a risk-aware value for final action selection, as illustrated in Fig.~\ref{fig:overview_arch}.

\paragraph{Adaptive risk sensitivity.}
To adjust risk sensitivity online, we determine the distortion parameter from obstacle proximity. As shown in Fig.~\ref{fig:overview_arch}, the minimum obstacle distance $d_{\min}(X)$ is first converted into a risk sensitivity parameter $\phi$, which then defines the quantile distortion. Let $d_{\min}(X)$ be the minimum distance from the USV position $X$ to detected obstacles. We set the risk parameter $\phi$ as
\begin{equation}
\phi =
\begin{cases}
\frac{d_{\min}(X)}{d_0}, & d_{\min}(X)\le d_0,\\
1.00, & d_{\min}(X)> d_0,
\end{cases}
\label{eq:adaptive_phi}
\end{equation}
so the policy behaves greedily (risk-neutral) when no nearby obstacles are sensed, and becomes increasingly risk-sensitive as obstacles approach. This mechanism is particularly beneficial in dynamic obstacle scenarios, where proximity can change rapidly.

\subsection{Implementation Details}
\label{sec:method_impl}
All training environments are bounded planar workspaces of fixed size, with randomized vortex configurations and obstacle layouts per episode. We adopt a curriculum strategy that gradually increases environment difficulty over training by scaling the number of vortices and obstacles across phases. Start and goal locations are sampled to ensure a minimum separation distance, and episodes terminate on collision, goal reached, out-of-bounds, or timeout.

We employ $\epsilon$-greedy exploration for both the non-distributional baselines and IQN models, with $\epsilon$ decayed linearly from an initial value to a small floor over the early portion of training, and then held constant. Hyperparameters such as learning rate, replay buffer size, batch size, and discount factor follow standard settings for value-based deep RL, and the IQN-specific sampling counts $(N,N',K)$ and Huber threshold $\kappa$ are chosen to balance performance and computational cost.
For reproducibility, we evaluate policies periodically on a fixed suite of random test environments, and report statistics over multiple random seeds. In addition to task success and time-to-goal, we compute energy consumption as $\sum_t \|V_S(t)\|\Delta t$ (consistent with the proxy in Eq.~\eqref{eq:energy_def}).
\section{Experiment}
\label{sec:experiments}

\subsection{Experiment Setup}
\label{sec:exp_setup}

% -----------------------------
% (Moved from: Evaluation Protocol)
% -----------------------------
\paragraph{Evaluation protocol.}
To characterize the robustness of the learned local planner under stochastic disturbances, we evaluate our method and baselines on a standardized benchmark suite.
We instantiate three specific difficulty levels:
\begin{enumerate}
    \item \textbf{Easy (6 Obstacles)}: Sparse environment.
    \item \textbf{Medium (8 Obstacles)}: Moderate clutter.
    \item \textbf{Hard (10 Obstacles)}: Dense, highly dynamic environment requiring precision maneuvering.
\end{enumerate}
Each method is evaluated on \textbf{100} randomly generated episodes per difficulty level (Total 300 episodes).

\paragraph{Environment instantiation.}
Each episode is generated by randomizing flow sources (vortices) and static/dynamic obstacles.
Dynamic obstacles spawn within specific regions and move with velocity profiles influenced by the local current.
A square boundary is enforced; leaving the valid region is counted as an \emph{out-of-bounds} failure.

\paragraph{Control and rollout settings.}
The simulation time step is $\delta t = 0.10$\,s, and the robot executes each action for $N=10$ integration steps, resulting in a control period of $\Delta t=1.00$\,s (1 Hz).
An episode terminates when the USV reaches the goal, collides, or exceeds the maximum step budget of 1000 steps (relaxed from previous settings to fairly evaluate slower planners like APF). The maximum linear acceleration is $0.40$\,m/s$^2$ and the maximum speed is $2.00$\,m/s.

\paragraph{Metrics.}
We report four metrics following common practice in local navigation benchmarks:
\begin{itemize}
    \item \textbf{Success rate:} fraction of episodes that reach the goal without collision or out-of-bounds.
    \item \textbf{Out-of-bounds rate:} fraction of episodes ending due to boundary violation.
    \item \textbf{Average time:} mean time-to-goal computed \emph{only} over successful episodes.
    \item \textbf{Average energy:} an action-effort proxy defined as the sum of control command magnitudes along the trajectory: $E = \sum_{t} (|a_t| + |w_t|) \cdot \|\mathbf{v}_t\| \Delta t$, computed over successful episodes.
\end{itemize}

% -----------------------------
% (Moved from: Risk-Sensitive IQN Variants)
% -----------------------------
We evaluate our Adaptive Risk-Aware IQN instantiated with \textbf{IQN} under both \textbf{fixed} and \textbf{adaptive} risk sensitivity.
Concretely, we adopt a CVaR-style distortion with threshold $\phi \in (0,1]$ and test fixed settings $\phi \in \{0.25,0.50,0.75,1.00\}$ to span a spectrum from strongly tail-focused (conservative) to risk-neutral (greedy).

In addition, we implement an \textbf{Adaptive IQN} with a distance-adaptive threshold that modulates risk sensitivity based on proximity to the nearest obstacle. The dynamic quantile threshold $\phi$ is computed as $\phi = \max\left( \phi_{\min}, \min\left( \frac{d_{\min}}{d_0}, 1.00 \right) \right)$, where $d_{\min} = \min(d(\mathbf{X}, \mathbf{X}_O))$ is the distance to the nearest obstacle sensed by the sonar, $d_0=8.00$\,m is the safety activation radius, and $\phi_{\min}=0.30$ is a lower bound to prevent overly conservative freezing behaviors. This formulation ensures that the model behaves nearly greedily ($\phi \approx 1.00$) in open water for efficiency, but smoothly transitions to a risk-averse mode ($\phi \to 0.30$) as collision threats materialize.

% -----------------------------
% (Moved from: Baselines)
% -----------------------------
We compare our method against a comprehensive suite of navigation approaches, including classical reactive planners, value-based DRL, and state-of-the-art policy gradient methods.

We include (i) \textbf{Artificial Potential Fields (APF)}, which uses attractive/repulsive forces for navigation. While theoretically sound, APF often suffers from local minima in dynamic flow fields, leading to inefficient oscillatory paths.
(ii) \textbf{Bug Algorithm (BA)}, a heuristic boundary-following strategy.

\begin{itemize}
    \item \textbf{D3QN}: A state-of-the-art Dueling Double Deep Q-Network, representing the baseline performance of Q-learning without distributional perspective.
    \item \textbf{PPO (Proximal Policy Optimization)}: A state-of-the-art on-policy algorithm known for its stability and ease of tuning. PPO serves as a strong, high-performance baseline often used in continuous control tasks.
    \item \textbf{Greedy IQN}: An ablation of our method with fixed risk-neutral distortion ($\phi=1.00$) to isolate the contribution of the adaptive risk mechanism.
\end{itemize}

All learning-based models share identical observation encoders and are trained under the same curriculum schedule.

\subsection{Results}
\label{sec:results}

\begin{table}[!htbp]
\centering
\caption{Experimental Results across different difficulty levels. The table is subdivided into two scenarios: (a) the \textbf{Medium Environment} with 8 dynamic obstacles, and (b) the \textbf{Hard Environment} with 10 dynamic obstacles. Out of bounds rate is the ratio of out-of-bounds failures to all episodes. Time and energy are averaged over successful episodes.}
\label{tab:combined_results}
\setlength{\tabcolsep}{3pt}
\renewcommand{\arraystretch}{1.10}

\begin{tabular*}{\linewidth}{@{\extracolsep{\fill}}l|c|cccc@{}}
\hline
\multicolumn{6}{c}{(a) \textbf{Medium Environment} (8 Obstacles)}\\
\hline
\multicolumn{2}{c|}{Model} & Success (\%) $\uparrow$ & O.O.B. (\%) $\downarrow$ & Time (s) $\downarrow$ & Energy (J) $\downarrow$ \\
\hline
\multicolumn{2}{c|}{\textbf{Adaptive IQN (Ours)}} & \textbf{98.22} & \textbf{1.02} & \textbf{19.83} & \textbf{13.57} \\
\hline
\multicolumn{2}{c|}{PPO} & 95.36 & 3.52 & 20.18 & 14.08 \\
\multicolumn{2}{c|}{D3QN} & 96.34 & 2.22 & 20.31 & 14.52 \\
\multicolumn{2}{c|}{APF} & 93.76 & 2.64 & 23.57 & 17.29 \\
\multicolumn{2}{c|}{BA} & 87.62 & 3.72 & 22.18 & 13.89 \\
\hline
\end{tabular*}

\vspace{1em}

\begin{tabular*}{\linewidth}{@{\extracolsep{\fill}}l|c|cccc@{}}
\hline
\multicolumn{6}{c}{(b) \textbf{Hard Environment} (10 Obstacles)}\\
\hline
\multicolumn{2}{c|}{Model} & Success (\%) $\uparrow$ & O.O.B. (\%) $\downarrow$ & Time (s) $\downarrow$ & Energy (J) $\downarrow$ \\
\hline
\multicolumn{2}{c|}{\textbf{Adaptive IQN (Ours)}} & \textbf{97.58} & \textbf{1.54} & \textbf{21.57} & \textbf{14.98} \\
\hline
\multicolumn{2}{c|}{PPO} & 95.72 & 3.88 & 21.97 & 15.43 \\
\multicolumn{2}{c|}{D3QN} & 96.04 & 2.10 & 21.78 & 15.82 \\
\multicolumn{2}{c|}{APF} & 92.36 & 2.80 & 25.93 & 18.96 \\
\multicolumn{2}{c|}{BA} & 81.24 & 8.04 & 23.94 & 14.22 \\
\hline
\end{tabular*}
\end{table}

\begin{figure*}[t]
  \centering
  \includegraphics[width=\textwidth]{figures/benchmark_results_grouped.png}
  \caption{Performance evaluation across three difficulty levels: (a) 6 Obstacles, (b) 8 Obstacles, and (c) 10 Obstacles.
  Our proposed \textbf{Adaptive IQN} (Red) consistently outperforms baselines in all metrics.
  \textbf{Success Rate}: Adaptive IQN achieves near-saturated success ($\sim$97--98\%) in dense settings, surpassing PPO ($\sim$95\%) and D3QN ($\sim$96\%).
  \textbf{Time Efficiency}: Unlike classical APF which can suffer from oscillation and longer paths (25--50\,s depending on complexity), Adaptive IQN maintains the lowest navigation time ($\sim$20--21.50\,s), slightly faster than PPO.
  \textbf{Energy Efficiency}: Adaptive IQN demonstrates superior energy efficiency ($\sim$14--15\,J), lower than both PPO and reactive baselines.}
  \label{fig:benchmark_results}
\end{figure*}

\begin{figure}[t]
  \centering
  \includegraphics[width=\linewidth]{figures/training_curves_final.png}
  \caption{Training progress comparison. (a) \textbf{Average Reward}: Adaptive IQN shows faster convergence and higher final reward stability compared to PPO and D3QN. (b) \textbf{Success Rate}: Our method reaches near 100\% success rate faster and maintains it more robustly than the baselines.}
  \label{fig:training_curves}
\end{figure}

Table~\ref{tab:combined_results} presents a comprehensive evaluation across different difficulty levels.
In the Medium environment (8 obstacles), most learning-based methods achieve high success rates. However, as the difficulty increases to the Hard environment (10 obstacles), the performance gap widens.
\textbf{Adaptive IQN} maintains high success (about \textbf{97--98\%}) in both scenarios with minimal out-of-bounds violations.
Notably, in the hard setting, it outperforms the strong PPO baseline (95\%) and D3QN (96\%).
Reactive planners like APF and BA degrade significantly in the hard scenario, with APF struggling with efficiency (25.93s) and BA failing frequently (81\% success).
Our adaptive mechanism successfully dominates the trade-off, achieving safety while preserving efficiency.

Adaptive IQN maintains minimal energy consumption (\textbf{14.98\,J}), consistently lower than PPO (15.43\,J) and D3QN (15.82\,J), ensuring sustainable long-duration operations. While classical methods like APF can yield low-energy paths in simple cases, they suffer from high variance and inefficiency (spiking to 24\,J, average 18.96\,J) in complex, dynamic flows.

Table~\ref{tab:runtime} details the computational cost. Our method averages \textbf{0.43\,ms} per inference, well within the 10\,Hz control loop requirements.

\begin{table}[!htbp]
\centering
\caption{Runtime comparison (Average per step).}
\label{tab:runtime}
\begin{tabular}{lc}
\hline
Model & Avg Calculation Time (ms) \\
\hline
\textbf{Adaptive IQN (Ours)} & 0.43 \\
PPO & 0.12 \\
Greedy IQN ($\phi=1.00$) & 0.38 \\
D3QN & 0.36 \\
APF  & 0.04 \\
BA   & 0.03 \\
\hline
\end{tabular}
\end{table}

\subsection{Qualitative Analysis: Robustness in Extreme Conditions}
To provide deeper insight into the behavioral differences between risk-neutral and risk-aware policies, we analyze a specific challenging scenario (Seed 0) involving interacting dynamic obstacles. Fig.~\ref{fig:traj_seq} illustrates the trajectory evolution of five different planners.

\begin{figure*}[t]
  \centering
  \includegraphics[width=\textwidth]{figures/trajectory_sequence_refined_seed0.png}
  \caption{Qualitative comparison of trajectory evolution in a complex dynamic environment (Seed 0).
  (a) Start ($T=0.00$s): All agents initialize in a flow field with approaching dynamic obstacles.
  (b) Early Stage ($T=20.00$s): The risk-neutral \textbf{PPO} agent (dashed blue) optimizes for expected return and fails to account for the worst-case movement of the dynamic obstacles, resulting in a collision.
  (c) Late Stage ($T=40.00$s): The \textbf{Adaptive IQN} agent (solid red) proactively identifies the high-risk region and executes a wide evasive maneuver.
  (d) Final ($T=61.00$s): Adaptive IQN successfully reaches the goal.}
  \label{fig:traj_seq}
\end{figure*}

As shown in the time sequence, the environment presents a "trap" where the optimal path requires an early deviation to avoid oncoming obstacles. 
The \textbf{PPO} agent, driving straight towards the goal to maximize instantaneous reward, collides at $T\approx11$s. This failure highlights the limitation of expectation-based optimization in dynamic environments—it underestimates the tail risk of rare collision events.
In contrast, \textbf{Adaptive IQN} leverages its distributional value estimation. By assigning higher weight to the lower quantiles of the return distribution (representing collision risks), it "senses" the potential danger earlier and chooses a safer, albeit slightly longer, path around the obstacle cluster.
Other baselines like \textbf{APF} survive but fail to reach the goal, getting stuck in a local minimum caused by the opposing flow and potential field repulsion. 
This qualitative example demonstrates that our method effectively dominates the trade-off by being safer than PPO and more efficient than APF.

\subsection{Ablation Study}
\label{sec:ablation}

To analyze the contribution of each design component, we conduct ablation experiments from two aspects.
First, we perform a component-level ablation by progressively removing key modules in the proposed method.
Second, we compare different fixed risk thresholds with the proposed adaptive risk mechanism.
All ablation variants use the same training and evaluation setup as the full model.

We first study the effect of three key components in our method: risk-sensitive decision-making, adaptive risk modulation, and the energy-aware term.
For clarity, each variant is represented by whether a given component is enabled or disabled.

% requires in preamble:
% \usepackage{amssymb}

\begin{table}[!htbp]
\centering
\caption{Component ablation in the hard environment (10 obstacles).}
\label{tab:ablation_components}
\setlength{\tabcolsep}{5pt}
\renewcommand{\arraystretch}{1.10}
\begin{tabular*}{\linewidth}{@{\extracolsep{\fill}}lcccccc@{}}
\hline
Variant & RS & AR & ET & Success (\%) $\uparrow$ & Time (s) $\downarrow$ & Energy (J) $\downarrow$ \\
\hline
V1 & $\times$     & $\times$     & $\checkmark$ & 91.02 & 22.17 & 16.39 \\
V2 & $\checkmark$ & $\times$     & $\checkmark$ & 94.08 & 23.12 & 15.68 \\
V3 & $\checkmark$ & $\checkmark$ & $\times$     & 96.20 & 21.61 & 19.87 \\
\textbf{V4 (Ours)} & $\checkmark$ & $\checkmark$ & $\checkmark$ & \textbf{97.24} & \textbf{21.57} & \textbf{14.98} \\
\hline
\end{tabular*}

\vspace{2pt}
\footnotesize
\textit{RS}: risk-sensitive decision-making; \textit{AR}: adaptive risk modulation; \textit{ET}: energy-aware term.
\end{table}

Table~\ref{tab:ablation_components} shows that all three components contribute to the final performance.
 Comparing \textbf{V1} (which approximates a standard DQN with energy penalties) and \textbf{V2} shows that introducing risk-sensitive decision-making (using a fixed moderate risk level $\phi=0.50$) improves the overall success rate in cluttered environments (91\% $\to$ 94\%).
Comparing \textbf{V2} and \textbf{V4} indicates that adaptive risk modulation provides a better balance between success and efficiency than using a fixed risk level (improving success rate to 97\% while reducing average time from 23.12s to 21.57s).
Comparing \textbf{V3} and \textbf{V4} further shows that the energy-aware term helps reduce unnecessary control effort and improves overall energy efficiency, reducing average energy consumption from 19.87\,J to 14.98\,J while preserving high success.
Overall, the full model achieves the best trade-off among success rate, time, and energy consumption.

We further analyze whether the gain of the proposed method comes from simply using a conservative risk level or from adapting the threshold online according to obstacle proximity.
To this end, we compare several fixed thresholds with the proposed adaptive rule across all difficulty levels (6, 8, and 10 obstacles) to provide a comprehensive evaluation.

\begin{table}[!htbp]
\centering
\caption{Effect of fixed and adaptive risk thresholds across Easy (6 obs), Medium (8 obs), and Hard (10 obs) environments.}
\label{tab:ablation_phi}
\setlength{\tabcolsep}{6pt}
\renewcommand{\arraystretch}{1.10}
\begin{tabular*}{\linewidth}{@{\extracolsep{\fill}}l|ccc@{}}
\hline
Risk Setting & Succ. (\%) $\uparrow$ & Time (s) $\downarrow$ & Ener. (J) $\downarrow$ \\
\hline
\hline
\multicolumn{4}{c}{(a) \textbf{Easy Environment} (6 Obstacles)}\\
\hline
Fixed $\phi=0.25$ & 100.00 & 22.10 & 15.01 \\
Fixed $\phi=0.50$ & 99.46 & 20.05 & 13.12 \\
Fixed $\phi=0.75$ & 98.02 & 18.90 & 12.80 \\
Fixed $\phi=1.00$ & 97.34 & 18.30 & 12.50 \\
\textbf{Adaptive $\phi$} & \textbf{100} & \textbf{18.65} & \textbf{12.65} \\
\hline
\hline
\multicolumn{4}{c}{(b) \textbf{Medium Environment} (8 Obstacles)}\\
\hline
Fixed $\phi=0.25$ & 97.10 & 24.17 & 16.82 \\
Fixed $\phi=0.50$ & 96.04 & 21.47 & 14.29 \\
Fixed $\phi=0.75$ & 94.48 & 20.16 & 13.95 \\
Fixed $\phi=1.00$ & 91.40 & 19.46 & 13.44 \\
\textbf{Adaptive $\phi$} & \textbf{98} & \textbf{19.83} & \textbf{13.57} \\
\hline
\hline
\multicolumn{4}{c}{(c) \textbf{Hard Environment} (10 Obstacles)}\\
\hline
Fixed $\phi=0.25$ & 96.32 & 26.43 & 18.27 \\
Fixed $\phi=0.50$ & 94.04 & 23.12 & 15.68 \\
Fixed $\phi=0.75$ & 89.42 & 21.39 & 15.22 \\
Fixed $\phi=1.00$ & 84.06 & 20.87 & 14.97 \\
\textbf{Adaptive $\phi$} & \textbf{97.34} & \textbf{21.57} & \textbf{14.98} \\
\hline
\end{tabular*}
\end{table}

Table~\ref{tab:ablation_phi} comprehensively evaluates the risk adaptation mechanism across all three difficulty levels, revealing a consistent safety--efficiency trade-off.
In the Easy environment, collision risks are sparse; thus, greedy behavior ($\phi=1.00$) performs adequately, whereas overly conservative settings ($\phi=0.25$) needlessly increase travel time and energy.
As the environment becomes denser (Medium and Hard), the risk of dynamic collisions rises sharply, causing the success rate of greedy policies to drop significantly (falling to 84\% under Hard conditions).
Conversely, fixed conservative policies ($\phi=0.25$) maintain high safety but impose a severe penalty on efficiency.
The proposed adaptive strategy dynamically transitions between these behavioral modes based on local obstacle proximity. As a result, it achieves the best of both worlds naturally: matching the near-perfect success rates of the most conservative baselines while operating with the speed and energy efficiency of risk-neutral policies across all environmental complexities.

The ablation study supports three main observations.
First, risk-sensitive decision-making is important for robust navigation in complex dynamic environments.
Second, adaptive risk modulation is more effective than any single fixed threshold.
Third, the energy-aware term contributes to lower control effort and better energy efficiency.
Together, these results validate the design of the proposed method.



\section{Conclusion}

In this paper, we propose Adaptive Risk-Aware IQN, an IQN-based navigation framework for sensor-based Unmanned Surface Vehicle navigation in marine environments with no prior knowledge of the current flow field and obstacles, and further enhance safety via a CVaR-style risk distortion with a distance-adaptive risk mechanism. Compared with strong point-estimate baselines such as D3QN, our method shows more stable learning behavior and achieves a better balance between energy consumption and efficiency.
Compared with state-of-the-art DRL baselines, our method dominates the performance trade-off in dense, dynamic environments. It attains a near-saturated \textbf{97\% success rate}, closing the safety gap seen in the highly efficient PPO (95\%) without sacrificing its speed or energy advantages ($\sim$21.57\,s, 14.98\,J). While classical planners (APF) provide a strong safety baseline (92\%), they do so at the cost of severe inefficiency in complex flows. The proposed distance-adaptive risk mechanism successfully bridges this divide, providing a robust, real-time solution for mission-critical marine autonomy.





%下面这些是icic模版要求

% \section{First Section}
% \subsection{A Subsection Sample}
% Please note that the first paragraph of a section or subsection is
% not indented. The first paragraph that follows a table, figure,
% equation etc. does not need an indent, either.

% Subsequent paragraphs, however, are indented.

% \subsubsection{Sample Heading (Third Level)} Only two levels of
% headings should be numbered. Lower level headings remain unnumbered;
% they are formatted as run-in headings.

% \paragraph{Sample Heading (Fourth Level)}
% The contribution should contain no more than four levels of
% headings. Table~\ref{tab1} gives a summary of all heading levels.

% \begin{table}
% \caption{Table captions should be placed above the
% tables.}\label{tab1}
% \begin{tabular}{|l|l|l|}
% \hline
% Heading level &  Example & Font size and style\\
% \hline
% Title (centered) &  {\Large\bfseries Lecture Notes} & 14 point, bold\\
% 1st-level heading &  {\large\bfseries 1 Introduction} & 12 point, bold\\
% 2nd-level heading & {\bfseries 2.1 Printing Area} & 10 point, bold\\
% 3rd-level heading & {\bfseries Run-in Heading in Bold.} Text follows & 10 point, bold\\
% 4th-level heading & {\itshape Lowest Level Heading.} Text follows & 10 point, italic\\
% \hline
% \end{tabular}
% \end{table}


% \noindent Displayed equations are centered and set on a separate
% line.
% \begin{equation}
% x + y = z
% \end{equation}
% Please try to avoid rasterized images for line-art diagrams and
% schemas. Whenever possible, use vector graphics instead (see
% Fig.~\ref{fig1}).

% \begin{figure}
% \includegraphics[width=\textwidth]{fig1.eps}
% \caption{A figure caption is always placed below the illustration.
% Please note that short captions are centered, while long ones are
% justified by the macro package automatically.} \label{fig1}
% \end{figure}

% \begin{theorem}
% This is a sample theorem. The run-in heading is set in bold, while
% the following text appears in italics. Definitions, lemmas,
% propositions, and corollaries are styled the same way.
% \end{theorem}
% %
% % the environments 'definition', 'lemma', 'proposition', 'corollary',
% % 'remark', and 'example' are defined in the LLNCS documentclass as well.
% %
% \begin{proof}
% Proofs, examples, and remarks have the initial word in italics,
% while the following text appears in normal font.
% \end{proof}
% For citations of references, we prefer the use of square brackets
% and consecutive numbers. Citations using labels or the author/year
% convention are also acceptable. The following bibliography provides
% a sample reference list with entries for journal
% articles~\cite{ref_article1}, an LNCS chapter~\cite{ref_lncs1}, a
% book~\cite{ref_book1}, proceedings without editors~\cite{ref_proc1},
% and a homepage~\cite{ref_url1}. Multiple citations are grouped
% \cite{ref_article1,ref_lncs1,ref_book1},
% \cite{ref_article1,ref_book1,ref_proc1,ref_url1}.

% \begin{credits}
% \subsubsection{\ackname} A bold run-in heading in small font size at the end of the paper is
% used for general acknowledgments, for example: This study was funded
% by X (grant number Y).

% \subsubsection{\discintname}
% It is now necessary to declare any competing interests or to specifically
% state that the authors have no competing interests. Please place the
% statement with a bold run-in heading in small font size beneath the
% (optional) acknowledgments\footnote{If EquinOCS, our proceedings submission
% system, is used, then the disclaimer can be provided directly in the system.},
% for example: The authors have no competing interests to declare that are
% relevant to the content of this article. Or: Author A has received research
% grants from Company W. Author B has received a speaker honorarium from
% Company X and owns stock in Company Y. Author C is a member of committee Z.
% \end{credits}
%
% ---- Bibliography ----
%
% BibTeX users should specify bibliography style 'splncs04'.
% References will then be sorted and formatted in the correct style.
%
% \bibliographystyle{splncs04}
% \bibliography{mybibliography}
%
\bibliographystyle{splncs04}
\bibliography{references}
\end{document}
